% GENERATED FILE. Edit canonical lessons, then run npm run generate:books.
% canonical-source-hash: fnv1a64:492a86e491068b83
% canonical-lessons: FR-C435-panneau, FR-C435-fleche, FR-C435-croisement, FR-C435-borne, FR-C435-angle

\chapter{Les panneaux}
\label{ch:fr-les-panneaux}
\hlchaptermodality{435}

\begin{chapteropening}
\textbf{By the end of this chapter:} \emph{I can follow signs: the sign, the arrow, the crossroads, the marker post and the corner.}

Complete the last lesson of chapter 435: I can follow signs: the sign, the arrow, the crossroads, the marker post and the corner.
\end{chapteropening}

\section[le panneau]{le panneau — the road sign; the panel}

\label{lesson:FR-C435-panneau}



\begin{quote}
Before the new one: say the French for the landmark, then the French for to mislay.
\end{quote}

\subsection*{You'll want to know: le panneau}
\textbf{le panneau} — \textquotedblleft{}the road sign; the panel\textquotedblright{}.

\begin{quote}
Le panneau indique le centre-ville à deux kilomètres.
\end{quote}

\begin{quote}
\emph{The sign shows the town centre two kilometres away.}
\end{quote}

Plural \emph{les panneaux}. \emph{Tomber dans le panneau} — to fall for it.

\begin{cousinweb}
\emph{panneau} comes from Old French \emph{panel} (EtymDB), from Vulgar Latin \emph{pannellus}, a small \emph{pannus} 'piece of cloth'.
\end{cousinweb}

\begin{grammarlens}[title={What you've built}]
One more word for telling, explaining and describing at B1.
\end{grammarlens}

\subsection*{Your turn}
\begin{itemize}
\raggedright
  \item \emph{Say it:} \emph{le panneau}
  \item \emph{Say it:} \emph{le panneau}, once more
  \item \emph{Say it:} a sentence of your own with \emph{le panneau}
  \item \emph{Recall:} say the French for the landmark, then the French for to mislay, then say \emph{le panneau} again
\end{itemize}

\begin{tcolorbox}[breakable,colback=teal!4,colframe=teal!35!black,title={Before you move on}]
What does le panneau mean? (\textquotedblleft{}The road sign; the panel\textquotedblright{}.) Say it once more.
\end{tcolorbox}
\section[la flèche]{la flèche — the arrow}

\label{lesson:FR-C435-fleche}



\begin{quote}
Before the new one: say the French for to mislay, then the French for the road sign.
\end{quote}

\subsection*{You'll want to know: la flèche}
\textbf{la flèche} — \textquotedblleft{}the arrow\textquotedblright{}.

\begin{quote}
Il suit les flèches jusqu'au parking.
\end{quote}

\begin{quote}
\emph{He follows the arrows to the car park.}
\end{quote}

\emph{Partir en flèche} — to shoot off.

\begin{cousinweb}
\emph{flèche} comes from Old French \emph{fleche}, which goes back to a Frankish word for 'arrow' (EtymDB).
\end{cousinweb}

\begin{grammarlens}[title={What you've built}]
One more word for telling, explaining and describing at B1.
\end{grammarlens}

\subsection*{Your turn}
\begin{itemize}
\raggedright
  \item \emph{Say it:} \emph{la flèche}
  \item \emph{Say it:} \emph{la flèche}, once more
  \item \emph{Say it:} a sentence of your own with \emph{la flèche}
  \item \emph{Recall:} say the French for to mislay, then the French for the road sign, then say \emph{la flèche} again
\end{itemize}

\begin{tcolorbox}[breakable,colback=teal!4,colframe=teal!35!black,title={Before you move on}]
What does la flèche mean? (\textquotedblleft{}The arrow\textquotedblright{}.) Say it once more.
\end{tcolorbox}
\section[le croisement]{le croisement — the crossroads, the junction}

\label{lesson:FR-C435-croisement}



\begin{quote}
Before the new one: say the French for the road sign, then the French for the arrow.
\end{quote}

\subsection*{You'll want to know: le croisement}
\textbf{le croisement} — \textquotedblleft{}the crossroads, the junction\textquotedblright{}.

\begin{quote}
Au croisement, nous tournons à droite.
\end{quote}

\begin{quote}
\emph{At the crossroads we turn right.}
\end{quote}

Close to \emph{le carrefour}.

\begin{cousinweb}
\emph{croisement} is built on \emph{croiser} 'to cross' with \emph{-ment}; \emph{croiser} comes from \emph{croix}, inherited from Latin \emph{crux} (Lewis \& Short). EtymDB does not list the noun.
\end{cousinweb}

\begin{grammarlens}[title={What you've built}]
One more word for telling, explaining and describing at B1.
\end{grammarlens}

\subsection*{Your turn}
\begin{itemize}
\raggedright
  \item \emph{Say it:} \emph{le croisement}
  \item \emph{Say it:} \emph{le croisement}, once more
  \item \emph{Say it:} a sentence of your own with \emph{le croisement}
  \item \emph{Recall:} say the French for the road sign, then the French for the arrow, then say \emph{le croisement} again
\end{itemize}

\begin{tcolorbox}[breakable,colback=teal!4,colframe=teal!35!black,title={Before you move on}]
What does le croisement mean? (\textquotedblleft{}The crossroads, the junction\textquotedblright{}.) Say it once more.
\end{tcolorbox}
\section[la borne]{la borne — the marker post, the boundary stone; the charging point}

\label{lesson:FR-C435-borne}



\begin{quote}
Before the new one: say the French for the arrow, then the French for the crossroads.
\end{quote}

\subsection*{You'll want to know: la borne}
\textbf{la borne} — \textquotedblleft{}the marker post, the boundary stone; the charging point\textquotedblright{}.

\begin{quote}
Une vieille borne au bord de la route indique le village.
\end{quote}

\begin{quote}
\emph{An old marker stone by the road points to the village.}
\end{quote}

\emph{Une borne kilométrique} — a kilometre stone; \emph{une borne de recharge} — a charging point.

\begin{cousinweb}
\emph{borne} comes from an Old French word found as \emph{bodne} and \emph{bontie} (EtymDB gives \emph{bontie}), which is usually traced to a Gaulish word for a boundary mark; the origin is not fully settled.
\end{cousinweb}

\begin{grammarlens}[title={What you've built}]
One more word for telling, explaining and describing at B1.
\end{grammarlens}

\subsection*{Your turn}
\begin{itemize}
\raggedright
  \item \emph{Say it:} \emph{la borne}
  \item \emph{Say it:} \emph{la borne}, once more
  \item \emph{Say it:} a sentence of your own with \emph{la borne}
  \item \emph{Recall:} say the French for the arrow, then the French for the crossroads, then say \emph{la borne} again
\end{itemize}

\begin{tcolorbox}[breakable,colback=teal!4,colframe=teal!35!black,title={Before you move on}]
What does la borne mean? (\textquotedblleft{}The marker post, the boundary stone; the charging point\textquotedblright{}.) Say it once more.
\end{tcolorbox}
\section[l'angle]{l'angle — the corner; the angle}

\label{lesson:FR-C435-angle}



\begin{quote}
Before the new one: say the French for the crossroads, then the French for the marker post.
\end{quote}

\subsection*{You'll want to know: l'angle}
\textbf{l'angle} — \textquotedblleft{}the corner; the angle\textquotedblright{}.

\begin{quote}
La pharmacie est à l'angle de la rue.
\end{quote}

\begin{quote}
\emph{The chemist's is on the corner of the street.}
\end{quote}

A masculine noun. \emph{À l'angle de} — on the corner of.

\begin{cousinweb}
\emph{angle} came through Old French from Latin \emph{angulus} 'corner, angle' (EtymDB).
\end{cousinweb}

\begin{grammarlens}[title={What you've built}]
One more word for telling, explaining and describing at B1.
\end{grammarlens}

\subsection*{Your turn}
\begin{itemize}
\raggedright
  \item \emph{Say it:} \emph{l'angle}
  \item \emph{Say it:} \emph{l'angle}, once more
  \item \emph{Say it:} a sentence of your own with \emph{l'angle}
  \item \emph{Recall:} say the French for the crossroads, then the French for the marker post, then say \emph{l'angle} again
\end{itemize}

\begin{tcolorbox}[breakable,colback=teal!4,colframe=teal!35!black,title={Before you move on}]
What does l'angle mean? (\textquotedblleft{}The corner; the angle\textquotedblright{}.) Say it once more.
\end{tcolorbox}
